\chapter{Methodology}
\label{cha:methodology}

The observed-data likelihood under the model of Chapter~\ref{cha:model-specification} is

\begin{equation}
    L(\bm{\theta})=\prod_{j=1}^{J}\prod_{i\in \mathcal{I}_{j}} p(\bm{y}_{ij}\mid \bm{\theta})=\prod_{j=1}^{J}\prod_{i\in \mathcal{I}_{j}} \int p(\bm{y}_{ij}\mid \bm{\lambda}_{j})p(\bm{\lambda}_{j}\mid \bm{\Sigma})\,d\bm{\lambda}_{j},
    \label{eq:log-lik-chap2}
\end{equation}
where the parameter space is $\bm{\theta}:=\{\bm{\mu}, \bm{\Phi}, \mathbf{B},\sigma^{2}\}$. This integral has no closed form in general because the Gaussian prior on $\bm{\lambda}_{j}$ is not conjugate to the multinomial likelihood, so we treat $\{\bm{\lambda}_{j}\}_{j=1}^{J}$ as missing data and then estimate $\bm{\theta}=\{\bm{\mu}, \bm{\Phi}, \mathbf{B},\sigma^{2}\}$ by EM algorithm. The standard approach in generalized linear mixed models is the Laplace approximation \citep{Breslow1993,Tierney1986}, which replaced the integral with a Gaussian approximation at the posterior mode.
% FIX: combined the two separate \citep's into one \citep{Breslow1993,Tierney1986}; two adjacent \citep's print as two separate parenthesized citations "(Breslow1993) (Tierney1986)" rather than "(Breslow1993; Tierney1986)".
To provide a more readable structure for the following sections, we introduce the basic notations and equalities for standard Laplace approximation.

Taking the logarithm of likelihood~\eqref{eq:log-lik-chap2} for each group $j$, we denote $h\left(\bm{\lambda}\right)$ as the logarithm of the integrand. The Laplace  method approximates the integral by forming a second-order Taylor expansion of $h\left(\bm{\lambda}\right)$ around its mode $\hat{\bm{\lambda}} = \arg\max_{\bm{\lambda}}h(\bm{\lambda})$.
% FIX: the argmax was written as "argmax ell(lambda)" but ell had not been defined yet -- h(lambda) was the name just introduced for this quantity. Changed to h(lambda) for consistency; if you did mean a distinct ell, please tell me how it relates to h and I will restate this properly. The rest of the chapter (Section~\ref{sec: EM-wb}) does use ell(lambda_j) for the specific group log-posterior, so if that is what you meant here, a forward pointer ("later specialized to ell(lambda_j) in Section...") would remove the ambiguity.
Let $\mathbf{H}$ be the negative Hessian matrix evaluated in the mode and $\mathbf{g}\left(\bm{\lambda}\right)$ be the gradient, then, 
\begin{equation}
    \mathbf{H}(\bm{\lambda}):=-\nabla^{2}h(\bm{\lambda})=\sum_{i\in\mathcal{I}_{j}}y_{i\cdot}\left(\text{diag}(\bm{p}_{i})-\bm{p}_{i}\bm{p}_{i}^{\top}\right)+\bm{\Sigma}^{-1},
    \label{eq: hessian-equality}
\end{equation}
\begin{equation}
    \mathbf{g}\left(\bm{\lambda}\right):=-\nabla h\left(\bm{\lambda}\right)=\sum_{i\in \mathcal{I}_{j}}\bigl(\bm{y}_{i}-y_{i\cdot}\bm{p}_{i}(\bm{\lambda})\bigr)-\bm{\Sigma}^{-1}\bm{\lambda}.
    \label{eq:gradient}
\end{equation}
For simplicity, we provide three definitions in advance.
% FIX: "For the simplicity" -> "For simplicity".
Define
\begin{equation*}
    d_{q}\left(\bm{\lambda}\right):=\sum_{i\in\mathcal{I}_{j}}y_{i\cdot}p_{iq}\left(\bm{\lambda}\right), \qquad \tilde{\mathbf{D}}\left(\bm{\lambda}\right):=\operatorname{diag}(\mathbf{d}\left(\bm{\lambda}\right))+\sigma^{-2}\mathbf{I}_{\text{Q}},
\end{equation*}
% FIX: the bracket around diag(...) was wrong -- it read diag(d(lambda) + sigma^{-2} I_Q), which adds a Q-vector to a Q x Q matrix inside diag(), a dimension mismatch. Moved sigma^{-2}I_Q outside diag(), matching the equivalent definition of D-tilde given later in this section (diag(d) + sigma^{-2}I_Q).
and the weight matrix $\mathbf{P}\left(\bm{\lambda}\right)\in \mathbb{R}^{I_{j}\times Q}$ with the $i$-th row $\sqrt{y_{i\cdot}}\bm{p}^{\top}_{i}\left(\bm{\lambda}\right)$, then the dense product can be rewritten as,
% FIX: P(lambda) was stated as in R^{Q x I_j}; since its i-th row is sqrt(y_i.)p_i^T (a 1xQ row vector) and there are I_j such rows, P must be I_j x Q, not Q x I_j. This also matches how P (and P^T) are used later: U(lambda):=[P(lambda)^T, B] is stated to be in R^{Q x (I_j+K)}, which requires P^T to be Q x I_j, i.e. P itself I_j x Q; and the block structure of the capacitance matrix Delta below only works out correctly if P is I_j x Q.
\begin{equation}
    \mathbf{P}^{\top}\mathbf{P}=\sum_{i\in\mathcal{I}_{j}}y_{i\cdot}\bm{p}_{i}\bm{p}_{i}^{\top},
    \label{eq:weight-prob}
\end{equation}
and finally define $\mathbf{M}_{K}:=\mathbf{I}_{K}+\sigma^{-2}\mathbf{B}^{\top}\mathbf{B}$.

\section{Woodbury Formula}
\label{sec:woodbury}

We introduce the Woodbury matrix identity, which we use throughout this section to invert matrices with a low-rank-plus-diagonal structure, such as $\bm{\Sigma}$ and the exact Hessian $\mathbf{H}$. For any invertible matrix $\mathbf{A}$ and matrices $\mathbf{U}, \mathbf{V}$ of compatible size, the identity reads
\begin{equation}
    \left(\mathbf{A} + \mathbf{U}\mathbf{V}^{\top}\right)^{-1} = \mathbf{A}^{-1} - \mathbf{A}^{-1}\mathbf{U}\left(\mathbf{I} + \mathbf{V}^{\top}\mathbf{A}^{-1}\mathbf{U}\right)^{-1}\mathbf{V}^{\top}\mathbf{A}^{-1}.
    \label{eq:standard-woodbury}
\end{equation}
In our setting the correction has rank $I_j+K$ rather than rank $K$ alone, since both the likelihood and the prior contribute a low-rank piece, so the Woodbury formula reduces the cost of inverting the Hessian from $\mathcal{O}(Q^3)$ to $\mathcal{O}\big(Q(I_j+K)^2+(I_j+K)^3\big)$, which is linear in $Q$ whenever the group size $I_j$ and the factor rank $K$ stay small relative to $Q$. By~\eqref{eq: hessian-equality} and \eqref{eq:weight-prob}, the Hessian can be written as the sum of a likelihood part and a prior part, where the likelihood part is
\begin{equation}
    \sum_{i\in\mathcal{I}_{j}}y_{i\cdot}\left(\text{diag}(\bm{p}_{i})-\bm{p}_{i}\bm{p}_{i}^{\top}\right)=\operatorname{diag}\left(\mathbf{d}\right)-\mathbf{P}^{\top}\mathbf{P},
    \label{eq:likelihood-part}
\end{equation}
with $d_q:=\sum_{i\in\mathcal{I}_j}y_{i\cdot}p_{iq}$. The prior part, transformed by the Woodbury formula, is
\begin{equation}
    \bm{\Sigma}^{-1} = \left(\mathbf{B}\mathbf{B}^{\top} + \sigma^2 \mathbf{I}_Q\right)^{-1} = \sigma^{-2}\mathbf{I}_Q - \sigma^{-4}\mathbf{B}\mathbf{M}_K^{-1}\mathbf{B}^{\top}.
    \label{eq:prior-part}
\end{equation}
Adding \eqref{eq:likelihood-part} and \eqref{eq:prior-part} together, we combine the diagonal terms into $\tilde{\mathbf{D}}:=\operatorname{diag}\left(\mathbf{d}\right)+\sigma^{-2}\mathbf{I}_{Q}$, while the two low-rank terms combine into
\begin{equation}
    \mathbf{P}^{\top}\mathbf{P} + \sigma^{-4}\mathbf{B}\mathbf{M}_K^{-1}\mathbf{B}^{\top}
    = \begin{bmatrix} \mathbf{P}^{\top} & \mathbf{B} \end{bmatrix}
    \begin{pmatrix} \mathbf{I}_{I_{j}} & \bm{0}\\ \bm{0} & \sigma^{-4}\mathbf{M}_K^{-1} \end{pmatrix}
    \begin{pmatrix} \mathbf{P} \\ \mathbf{B}^{\top} \end{pmatrix},
    \label{eq:combined}
\end{equation}
i.e., $\mathbf{U}\mathbf{C}\mathbf{U}^{\top}$. This gives \eqref{eq:hessian-decomp}. Positive semi-definiteness follows from $\mathbf{C}\succ 0$, and the rank does not exceed $I_{j}+K$, the number of columns of $\mathbf{U}$. Two points are worth noting here. First, the cross term $-\bm{p}_{i}\bm{p}_{i}^{\top}$ in the likelihood is exactly what makes $\mathbf{H}$ dense. Its row rank is only 1, and combined over $I_{j}$ observations the rank is at most $I_{j}$. Second, the low-rank part of the prior originates from the factor structure itself, with rank $K$. The two happen to fold into the same $\mathbf{U}$, and this is the entire source of the speed-up in this section. The "denseness" of the data and the "low-rankness" of the prior are, here, the same thing.

Thus, for any $\bm{\lambda}\in\mathbb{R}^{Q}$, define
\begin{equation}
    \mathbf{U}\left(\bm{\lambda}\right):=\left[\mathbf{P}\left(\bm{\lambda}\right)^{\top} \quad \mathbf{B}\right]\in \mathbb{R}^{Q\times (I_{j}+K)}, \qquad \mathbf{C}:=\begin{pmatrix}
    \mathbf{I}_{I_{j}}
    &
    \bm{0}
    \\
    \bm{0}
    &
    \sigma^{-4}\mathbf{M}_{K}^{-1}
    \end{pmatrix}.
    \label{eq:wb-decomp}
\end{equation}
Then the exact Hessian structure satisfies
\begin{equation}
    \mathbf{H}=\tilde{\mathbf{D}}-\mathbf{U}\mathbf{C}\mathbf{U}^{\top}.
    \label{eq:hessian-decomp}
\end{equation}
Applying \eqref{eq:standard-woodbury} to $\mathbf H=\tilde{\mathbf D}+\mathbf U(-\mathbf C)\mathbf U^\top$ gives the Woodbury inverse
\begin{equation}
    \mathbf{H}^{-1}=\left(\tilde{\mathbf{D}}-\mathbf{U}\mathbf{C}\mathbf{U}^{\top}\right)^{-1} =\tilde{\mathbf{D}}^{-1} + \tilde{\mathbf{D}}^{-1}\mathbf{U}\left(\mathbf{C}^{-1} - \mathbf{U}^{\top}\tilde{\mathbf{D}}^{-1}\mathbf{U}\right)^{-1}\mathbf{U}^{\top}\tilde{\mathbf{D}}^{-1}.
    \label{eq:hessian-decomp-wb}
\end{equation}

Let
\begin{equation}
    \Delta\left(\bm{\lambda}\right):= \mathbf{C}^{-1} - \mathbf{U}^{\top}\tilde{\mathbf{D}}^{-1}\mathbf{U}
= \begin{pmatrix}
\mathbf{I}_{I_{j}} - \mathbf{P}\tilde{\mathbf{D}}^{-1}\mathbf{P}^{\top} & -\mathbf{P}\tilde{\mathbf{D}}^{-1}\mathbf{B} \\
-\mathbf{B}^{\top}\tilde{\mathbf{D}}^{-1}\mathbf{P}^{\top} & \sigma^4\mathbf{M}_K - \mathbf{B}^{\top}\tilde{\mathbf{D}}^{-1}\mathbf{B}
\end{pmatrix}
    \label{eq:capacitance}
\end{equation}
denote the capacitance matrix, of the smaller size $(I_{j}+K)\times (I_{j}+K)$ rather than the Hessian's $Q\times Q$, so that $\mathbf{H}^{-1}=\tilde{\mathbf{D}}^{-1}+\tilde{\mathbf{D}}^{-1}\mathbf{U}\Delta^{-1}\mathbf{U}^{\top}\tilde{\mathbf{D}}^{-1}$. For the inverse of $\mathbf H$ to exist and for its Cholesky factorization to be valid, $\Delta$ must be positive definite. This is not an extra condition to check. $\tilde{\mathbf D}\succ 0$ trivially, and the log-posterior is strictly concave at the mode, so $\mathbf H\succ 0$. The matrix determinant lemma applied to \eqref{eq:hessian-decomp-wb} then forces $\Delta\succ 0$ as well.
% FLAG (not changed): a positive determinant alone does not imply positive definiteness -- an even number of negative eigenvalues also gives a positive determinant. The matrix determinant lemma by itself only shows det(Delta) has the same sign pattern implied by det(H), det(Dtilde), det(C); it does not rule out an even number of negative eigenvalues in Delta. Earlier in this conversation we worked through a proper proof of this positive-definiteness claim (your Theorem 4, via the Schur complement of a larger block matrix P = [[Dtilde, U],[U^T, C^{-1}]]), which is a genuinely different and more complete argument than "matrix determinant lemma forces Delta > 0". If that theorem is proved in full elsewhere in the thesis, I'd suggest citing it here instead of re-deriving a shortcut version that doesn't quite close the gap; if this paragraph is meant to BE the proof, it needs the Schur-complement step added.

Putting the pieces together, computing $\mathbf{H}^{-1}$, or a solve $\mathbf{H}^{-1}\bm v$, never requires factorizing a $Q\times Q$ matrix. We only need $\tilde{\mathbf D}^{-1}$, which is diagonal, the $(I_j+K)\times(I_j+K)$ matrix $\Delta(\bm\lambda)$, and a Cholesky factorization of $\Delta$. Since $I_j$ and $K$ are both small and fixed while $Q$ grows, this gives the $\mathcal O\big(Q(I_j+K)^2+(I_j+K)^3\big)$ cost stated at the start of this section, in place of the dense $\mathcal O(Q^3)$ inversion.

\section{Laplace-EM Algorithm via Woodbury}
\label{sec: EM-wb}

In this section, we construct a Laplace EM algorithm to estimate the parameters in the FA-DMR. The E-step relies on a Laplace approximation to compute the approximate models of the random effects. \citet{Rubin1982} established the Rubin-Thayer algorithm for the extraction of $\mathbf{B}$ and $\mathbf{D}$ from $\bm{\Sigma}$. Under our model assumption, the idiosyncratic matrix is determined by $\sigma^{2}$ rather than a diagonal matrix with different entries.
% FIX: "the idiosyncratic matrix determined by sigma^2" -> "...is determined by..." (missing verb).
Thus, we introduce one classical result in probabilistic principal component analysis (PPCA) by \citet{Tipping1999}, who had shown that,
% FIX: "principle component analysis" -> "principal component analysis".
\begin{equation}
    \mathbf{B}=\mathbf{U}_{K}\left(\bm{\Lambda}_{K}-\sigma^{2}\mathbf{I}\right)^{1/2}\mathbf{R}.
    \label{eq:tipping-ppca_b}
\end{equation}
For the estimated $\mathbf{B}$, $\sigma^{2}$ satisfies,
\begin{equation}
    \sigma^{2}=\frac{1}{Q-K}\sum_{n=K+1}^{Q}\lambda_{n},
    \label{eq:tipping-ppca_sigma}
\end{equation}
where the $K$ column vectors in the $Q\times K$ matrix $\mathbf{U}_{K}$ are the principal eigenvectors of the sample covariance matrix of the group index, with corresponding eigenvalues $\lambda_{1},\cdots,\lambda_{K}$ in the $K\times K$ diagonal matrix $\bm{\Lambda}_{K}$, and $\mathbf{R}$ is an arbitrary $K\times K$ orthogonal rotation matrix.
% FLAG (not changed): "the sample covariance matrix of the group index" is unclear to me -- an index is not itself a random variable with a covariance. Do you mean the sample covariance of the group-level random effects (i.e. the matrix Sigma-hat-obs from Equation~\eqref{eq:sigma-obs}, or S from the FH-EM section)? Please confirm the exact object and I will name it precisely.
In our model settings, the PPCA decomposition should be done in $(Q-1)$ dimensions, since the gauge constraint $\mathbf{1}^{\top}\mathbf{B}=\bm{0}$ gives one eigen-value to be exact 0. The transformed matrix $\bm{\Pi}$ for loading matrix $\mathbf{B}$ should be applied here for corresponding to the $(Q-1)\times (Q-1)$ dimensions as well. For the covariance matrix $\mathbf{S}$, the projection onto $\mathbf{1}^{\perp}$ is given by $\mathbf{S}_{\perp}:=\mathbf{V}^{\top}\mathbf{S}\mathbf{V}\in \mathbb{R}^{(Q-1)\times (Q-1)}$. Then we decompose the eigen-vector and find $Q-1$ eigen-values $\lambda_{1}\geq \lambda_{2}\geq \cdots\geq \lambda_{Q-1}$ and eigen-vector $\tilde{\mathbf{U}}\in \mathbb{R}^{(Q-1)\times (Q-1)}$. Let $\Lambda_{K}=\operatorname{diag}\left(\lambda_{1},\lambda_{2}, \cdots,\lambda_{K}\right)$, the new update formula for our $\sigma^{2}$ is given by,
\begin{equation}
    \hat{\sigma}^{2}=\frac{1}{Q-1-K}\sum_{n=K+1}^{Q-1}\lambda_{n},
    \label{eq:tipping-ppca_sigma_deflate}
\end{equation}
and corresponding $\hat{\mathbf{B}}$ calculated by the product of classical decomposition from $Q-1$ space and transform projection $\mathbf{V}$,
\begin{equation}
    \hat{\mathbf{B}}=\mathbf{V}\tilde{\mathbf{U}}_{\mathrm{K}}\left(\bm{\Lambda}_{\mathrm{K}}-\hat{\sigma}^{2}\mathbf{I}_\mathrm{K}\right)^{1/2}\mathbf{R}.
    \label{eq:tipping-ppca_b_deflate}
\end{equation}
The M-step updates model parameters via the Rubin-Thayer algorithm \citep{Rubin1982}. For the high dimensional data with large group number $J$ and category number $Q$, we introduce the Woodbury equality to reduce the calculation complexity and the runtime, without changing the estimate itself. Since Woodbury is an exact algebraic identity applied to the same multinomial-model quantities, the estimation accuracy is identical to the non-accelerated computation; only the runtime is reduced. Thus, we report results from the fast, Woodbury-accelerated algorithm, rather than the standard (non-accelerated) Laplace-EM algorithm, purely to reduce runtime.

\paragraph*{E-step: Estimate the Random Effect via LA.}
For each group $j$, the log-posterior of $\bm{\lambda}_{j}$ is given by:
\begin{equation}
    \ell(\bm{\lambda}_{j})=\sum_{i\in \mathcal{I}_{j}}\Bigl[\sum_{q}y_{ijq}\eta_{ijq}(\bm{\lambda}_{j})-y_{ij\cdot}\log \sum_{q}\exp\bigl(\eta_{ijq}(\bm{\lambda}_{j})\bigr)\Bigr]-\frac{1}{2}\bm{\lambda}_{j}^{\top}\bm{\Sigma}^{-1}\bm{\lambda}_{j},
    \label{eq: estep-loglik}
\end{equation}
up to an additive constant independent of $\bm{\lambda}_{j}$, where 
\begin{equation}
    \eta_{ijq}(\bm{\lambda}_{j})=\mu_{q}+\bm{v}_{ij}^{\top}\bm{\phi}_{q}+\lambda_{jq},
    \label{eq: individualcovariate}
\end{equation}
is evaluated separately for each $i\in \mathcal{I}_{j}$ to reflect the individual-specific covariate structure.
% FIX: "...gives the individual-specific covariate structure" -> "...to reflect the individual-specific covariate structure" (the original did not connect grammatically to the previous clause).
The gradient and Hessian of \eqref{eq: estep-loglik} are~\eqref{eq: hessian-equality} and \eqref{eq:gradient}. Since $\operatorname{diag}(\bm p_i)-\bm p_i\bm p_i^\top\succeq 0$ and $\bm{\Sigma}^{-1}\succ 0$, we have $\mathbf{H}(\bm{\lambda}_{j})\succ 0$ everywhere, so $-\nabla^2\ell(\bm\lambda_j) = \mathbf H(\bm\lambda_j) \succ 0$, i.e.\ $\ell(\bm{\lambda}_{j})$ is strictly concave, and has a unique maximizer,
% FIX (part of the H-sign cluster -- see my note above): the original wrote "we have H(lambda_j) prec 0 everywhere, so ell is strictly concave". Given H is defined as the NEGATIVE Hessian, and is shown to be PD here (sum of a PSD term and a PD term), H should be H > 0, not H < 0; concavity of ell then follows because -nabla^2 ell = H > 0 means nabla^2 ell < 0. I also swapped the undefined symbol "P_ij" for the explicit expression diag(p_i)-p_i p_i^T it evidently refers to, since P was already used earlier in this chapter for a different object (the stacked I_j x Q weight matrix) -- reusing P for this per-observation term risks confusion between the two. Please check this whole passage against your code.
\begin{equation}
     \hat{\bm{\lambda}}_{j}
     = \arg\max_{\bm{\lambda}\in \mathbb{R}^{Q}}
     \ell(\bm{\lambda}_{j}),
     \label{eq:maximizer}
\end{equation}
% FIX: changed "ell_j(lambda_j)" to "ell(lambda_j)" to match the notation ell(lambda_j) used in Equation~\eqref{eq: estep-loglik} (no subscript j on ell there).
which can be found by Newton's method.
% FLAG (not changed): the next sentence, "Strict concavity guarantees convergence to lambda-hat_j from any starting point", is a strong claim -- plain Newton's method (without a globalization step) is not guaranteed to converge globally from strict concavity alone; it typically needs a line search or trust region, which is exactly what the Armijo backtracking mentioned a few sentences later provides. As currently ordered, the strong global-convergence claim appears before backtracking is introduced, and sits oddly next to the later admission that "we observed severe instability" that motivated adding backtracking in the first place. Consider moving this claim to after backtracking is introduced, or weakening it to a local-convergence statement with global convergence attributed to the backtracking step.
Strict concavity guarantees convergence to $\hat{\bm{\lambda}}_{j}$ from any starting point. The Laplace approximation then replaces the intractable posterior with the Gaussian matched to the model and curvature at,
\begin{equation}
    \bm{\lambda}_{j}^{(t+1)}=
    \bm{\lambda}_{j}^{(t)}
    +s^{(t)}\mathbf{H}\bigl(\bm{\lambda}_{j}^{(t)}
    \bigr)^{-1}\nabla_{\bm{\lambda}_{j}}
    \ell \bigl(\bm{\lambda}_{j}^{(t)}\bigr),
    \label{eq:newton-method}
\end{equation}
% FIX (part of the H-sign cluster): changed the minus sign to a plus sign. With H defined as the negative Hessian (H = -grad^2 ell), the Newton ascent step for maximizing ell is lambda + s*H^{-1}*grad(ell), not lambda - s*H^{-1}*grad(ell); as written with a minus sign this moves in a descent direction. Please check against your code -- if the code uses a minus sign, either the code has the same bug, or H is defined with the opposite sign convention there and this write-up's H := -grad^2 h needs to change instead.
The Laplace approximation then yields,
\begin{equation}
    p(\bm{\lambda}_{j}\mid \bm{y}_{ij}) \approx \mathcal{N}\bigl(\hat{\bm{\lambda}}_{j}, \hat{\mathbf{S}}_{j}\bigr), \quad \hat{\mathbf{S}}_{j}=\mathbf{H}(\hat{\bm{\lambda}_{j}})^{-1}.
    \label{eq:surrogate-eq}
\end{equation}
% FIX (part of the H-sign cluster): removed the extra negation. H is already the negative Hessian (already positive definite at the mode), so the Laplace posterior covariance is H^{-1} directly; (-H)^{-1} would be negative definite, which cannot be a covariance matrix.
In E-step, we also designed the improvement using Multinomial-Poisson Equivalence in the calculation of Hessian before, where the first term with complexity $\mathcal{O}\left(Q^{3}\right)$ in~\eqref{eq: hessian-equality} was removed.
% FIX: "was been removed" -> "was removed".
However, we observed severe instability on the value of log-likelihood, leading to a challenging issue of non-monotonicity in optimization.
% FIX: "leading challenging issue of non-monotonicity and optimization" -> "leading to a challenging issue of non-monotonicity in optimization".
Meanwhile, the Heywood problem arises. We may encounter a severe Heywood problem, where the estimated error covariance matrix collapses during the iterations, with the final estimates degenerating toward zero. Therefore, we use Newton's method with Armijo backtracking search with step size $s^{(t)}\in (0,1]$ for the estimation of $\bm{\lambda}$. 

\paragraph*{M-step: Maximization and Estimation of Fixed Effects.}
With $\hat{\bm{\lambda}}_{j}$ and $\hat{\mathbf{S}}_{j}$ estimated, we solve the maximization problem independently, 
\begin{equation}
    (\hat{\mu}_{q}, \hat{\bm{\phi}}_{q})=\arg \max_{\mu_{q}, \bm{\phi}_{q}}\sum_{i}\bigl[y_{ijq}(\mu_{q}+\bm{v}_{ij}^{\top}\bm{\phi}_{q})-\exp (\mu_{q}+\bm{v}_{ij}^{\top}\bm{\phi}_{q}+\hat{\lambda}_{jq}+\xi_{ij})\bigr].
\end{equation}
This is a Poisson log-linear regression with offset $\lambda_{jq}+\xi_{ij}$, solved by iteratively reweighted least squares. The $Q$ problems are independent and may be parallelized. The EM sufficient statistic for $\bm{\Sigma}$ follows directly from the Laplace output:
\begin{equation}
    \hat{\bm{\Sigma}}_{\text{obs}}=\frac{1}{J}\sum_{j}\mathbb{E}[\bm{\lambda}_{j} \bm{\lambda}_{j}^{\top}\mid \bm{y}_{ij}]\approx \frac{1}{J}\sum_{j}(\hat{\bm{\lambda}}_{j} \hat{\bm{\lambda}}_{j}^{\top}+\hat{\mathbf{S}}_{j}).
    \label{eq:sigma-obs}
\end{equation}
Critically, $\hat{\bm{\Sigma}}_{\text{obs}}$ is fully determined by current Laplace output $(\hat{\bm{\lambda}}_{j}, \hat{\mathbf{S}}_{j})$ and does not need to be re-evaluated during the Rubin-Thayer iterations, thus eliminating the circular dependency inherent in factor-space formulations.
Given $\hat{\bm{\Sigma}}_{\text{obs}}$, we decompose $\bm{\Sigma}=\mathbf{BB}^{\top}+\sigma^2\mathbf{I}_{Q}$ via Equations~\eqref{eq:tipping-ppca_b}-\eqref{eq:tipping-ppca_sigma} \citep{Tipping1999}. The PLT constraint is enforced post-hoc via a QR decomposition of the top $K$ rows of $\hat{\mathbf{B}}_{\text{new}}$, followed by sign correction to ensure positive diagonal entries and zeroing of the upper triangle. Applying PLT only after the Rubin–Thayer sub-iterations converge, rather than at each inner step, avoids the discontinuities that would otherwise disrupt convergence.

Given $\hat{\mathbf{B}}$ and $\hat{\sigma}^2$, the factor scores are recovered by the Bayes linear estimator \citep{Gelman2013}:
\begin{equation}
\hat{\mathbf{f}}_j = \left( \mathbf{I}_K + \frac{1}{\hat{\sigma}^2} \hat{\mathbf{B}}^\top \hat{\mathbf{B}} \right)^{-1} \frac{1}{\hat{\sigma}^2} \hat{\mathbf{B}}^\top \hat{\boldsymbol{\lambda}}_j ,
\label{eq: Bayeslinearestimator}
\end{equation}
requiring only $K\times K$ inversion.  

\subsection{Limitations of the Standard EM-Algorithm \& Solutions}

Although the runtime of the Woodbury-EM algorithm reduces substantially compared to the standard EM algorithm and existing R packages such as \texttt{glmmTMB}, its accuracy is a separate question, because our work so far has focused only on acceleration, not on optimization.
% FIX: reworded "the accuracy is challenged, because our previous work all based on the acceleration rather than optimization" for grammar and clarity, keeping the same claim.
Small-scale tests revealed an interesting pattern: the estimation error of the loading matrix $\mathbf{B}$ accumulates through iterations, a consequence of how the EM algorithm works, resulting in a final estimate of $\mathbf{B}$ that is slightly less accurate than \texttt{glmmTMB} and \texttt{gllvm}.
% FIX: reworded "There are some facts that interesting when we did the small-scale tests, that... resulting the final estimation of B in slightly less accuracy than..." for grammar.
Meanwhile, the estimation of $\mathbf{B}$ affects $\sigma^{2}$ in this iteration, according to Equations~\eqref{eq:tipping-ppca_b}--\eqref{eq:tipping-ppca_sigma}; we refer to this as the coupling of the two steps.
% FIX: split the run-on sentence and changed \ref to \eqref.

For more details, we come back to the prior in E-step $\bm{\Sigma}^{-1}=\left(\hat{\mathbf{B}}\hat{\mathbf{B}}^{\top}+\sigma^{2}\mathbf{I}\right)^{-1}$, which affected by $\hat{\mathbf{B}}^{(t)}$ itself in $t$-th iteration. In practice, this point leads to a self-confirming feedback loop. Along the directions that have already been identified as important by $\hat{\mathbf{B}}$, the prior variance is large and there is less shrinkage, while in the orthogonal directions, the prior variance is only $\sigma^{2}$ and there is more shrinkage. Therefore, $\hat{\bm{\lambda}}_{j}$ mainly spreads along the directions that have already been identified by $\hat{\mathbf{B}}$, and the M-step then uses these $\hat{\lambda}_j$ to estimate $\hat{\mathbf{B}}$ back to the same directions. The systematic localization experiments, including two comparison groups with $\bm{\mu},\bm{\phi}$ fixed and $\mathbf{B},\sigma^2$ varying, and with $\mathbf{B},\sigma^2$ fixed and $\bm{\mu},\bm{\phi}$ varying, precisely locate the source of the problem in this mechanism, rather than the estimation of $\bm{\mu},\bm{\phi}$ or the Rubin-Thayer method itself.

In the following section, we arrange three methods of correction and compare the results in terms of accuracy, runtime, and robustness.
% FIX: "of both accuracy, runtime and robustness" -> "in terms of accuracy, runtime, and robustness" ("both" does not fit a list of three).

\section{Fay-Herriot EM Algorithm (FH-EM): Maximum Likelihood Path}
\label{sec:fay-herriot-em}

\paragraph*{Model Assumption.} The Fay-Herriot marginal likelihood in Equation~\eqref{eq:fh-marginal-likelihood} borrows its structure directly from the classical small-area model~\citep{Herriot1979, Pfeffermann2011}: a direct estimate observed with a known sampling covariance, linked through a common between-group Gaussian distribution. Transplanting this structure from survey-based small-area estimation to our setting requires several assumptions that do not hold exactly here, and we state them explicitly rather than leave them implicit in Equation~\eqref{eq:observation_model}.
% FLAG (not changed): the bib key here is "Herriot1979" (second author only); the founding paper is Fay & Herriot (1979). If this is your standing convention for this key that's fine, just flagging in case it was meant to be FayHerriot1979 or similar.

\begin{enumerate}
\item[(A1)] \emph{$\mathbf F_j$ is treated as known, though it is not.} In the classical model, the sampling covariance $\psi_i$ is a design quantity, fixed and known prior to estimation. Here $\mathbf F_j$ is the multinomial Fisher information evaluated at $\hat{\bm\lambda}_j^{\text{iso}}$, the mode of an auxiliary isotropic-prior Laplace approximation (Section~\ref{subsubsec:isotropic-prior}), and is therefore itself an estimated quantity that depends on the current global parameters $(\hat{\mathbf{B}},\hat\sigma^2)$ through $\tau^2$. Equation~\eqref{eq:observation_model} treats $\mathbf F_j$ as fixed and non-stochastic within one EM round; this is the same plug-in approximation already flagged for the marginal likelihood in Equation~\eqref{eq:gaussian-mix}, and it propagates into the E-step and M-step derivations exactly the same way. It is not an additional error introduced by the EM reformulation; it is the same approximation, made explicit at the point where it is used rather than only at the point where its consequence (the Gaussian-mixture correction) is visible.

\item[(A2)] \emph{Validity regime of the Gaussian approximation.} $\tilde{\bm\lambda}_j\mid\bm\lambda_j\sim\mathcal N(\bm 0,\mathbf F_j^{-1})$ relies on the asymptotic normality of the Laplace-approximated MLE. This is accurate when the per-category information within a group is large, and degrades as the average count per category shrinks, independent of $J$ or $Q$ themselves --- the same information-per-category condition identified in the debias failure analysis (Section~\ref{subsubsec:debias}). We do not assume this approximation is exact for any finite sample; we assume it is adequate in the regime our simulations and asymptotic theory (Chapter~\ref{cha:asymptotics}) target.

% FLAG (not changed): (A3) is missing -- the list jumps from (A2) directly to (A4). Please check whether an assumption was deleted without renumbering the rest, or whether the numbering itself just needs to close the gap (A1-A6).

\item[(A4)] \emph{Identifiability along the singular direction.} $\mathbf F_j$ is exactly singular along the all-ones direction, since a constant shift to all categories leaves the softmax probabilities unchanged. $\mathbf F_j+\bm\Sigma^{-1}$ is well-posed without an additional ridge because $\bm\Sigma^{-1}\succ0$ regularizes this direction on its own; however, $\tilde{\bm\lambda}_j$ itself is only defined up to whatever convention fixes this direction, and that convention is inherited from wherever $\bm\lambda_j$'s own identifiability constraint is imposed, not introduced fresh here.

\item[(A5)] \emph{Scope of the ascent-property claim.} The genuine $\arg\max$ structure and the ascent property established in this section apply to one inner update of $\bm\Sigma$ for a fixed $(\tilde{\bm\lambda}_j,\mathbf F_j)$ pair. Because $\mathbf F_j$ is itself re-estimated every outer round from the evolving isotropic-prior stage, as noted in (A2), the full nested procedure is a generalized EM in the sense that the auxiliary quantities being conditioned on change between rounds. We therefore claim the ascent property only for the inner FH-EM sub-problem, exactly as stated; we do not claim that a single fixed objective function is monotonically increased across outer rounds of the whole isotropic-plus-FH-EM procedure.

\item[(A6)] \emph{What the positive-definiteness guarantee does and does not cover.} $\mathbf S$ in Equation~\eqref{eq:fhem_mstep} is a sum of positive semi-definite terms and therefore cannot produce the negative-eigenvalue collapse that afflicts the debias estimator's moment subtraction (Section~\ref{subsubsec:debias}). This rules out that specific failure mode by construction. It does not imply that $\hat\sigma^2$ can never reach the numerical floor: if the data genuinely favor $\sigma^2\to0$ in the ordinary maximum-likelihood sense, FH-EM has no more immunity to this classical boundary solution than any other exact likelihood-based method, and this is a separate phenomenon from the one FH-EM is designed to correct.

\item[(A7)] \emph{Local, not global, optimum.} As with any EM algorithm applied to a factor-analytic latent-variable model, the likelihood surface need not be unimodal, and the fixed point reached depends on initialization. We do not claim FH-EM recovers the global maximum of Equation~\eqref{eq:fh-marginal-likelihood}; we claim that, given a fixed starting point, each step does not decrease the Q-function in Equation~\eqref{eq:fhem-Q-function}.
\end{enumerate}

Equation~\eqref{eq:debias-Sigma} is a moment matching formula rather than the $\arg\max$ of any objective function. Therefore, the ascent property of the standard EM algorithm does not apply to this step from the beginning. Its validity comes from a separate argument that the moment estimator is consistent, rather than from the EM theory itself. The collapse is a direct consequence of this point. If we want to add a boundary correction to this step to prevent the degeneration, similar to the idea of adjusted ML, there is no natural objective function to modify. It is not maximizing any objective function, so we cannot simply say that a penalty term is added to the objective function.

The first stage of the debias method, which uses the Laplace approximation with an isotropic prior, remains unchanged and still produces $(\tilde{\bm{\lambda}}_j,\mathbf{F}_j)$. In the second stage, we replace the moment matching step with a genuine $\arg\max$ problem. We define $\tilde{\bm{\lambda}}_{j}-\bm{\lambda}_{j}=\bm{\epsilon}_j$ and have the conditional distribution,

\begin{equation} 
    \bm{\epsilon}_{j}\mid \bm{\lambda}_{j} \sim \mathcal{N}\left(\mathbf{0}, \mathbf{F}_j^{-1}\right). 
    \label{eq:observation_model} 
\end{equation} 
Together with $\bm{\lambda}_{j}\sim \mathcal{N}\left(\bm{0},\bm{\Sigma}\right)$, we integrate on $\bm{\lambda}_{j}$ marginally. The density of $\tilde{\bm{\lambda}}_{j}$ is given by $p\left(\bm{\lambda}_{j},\tilde{\bm{\lambda}}_{j}\right)=p\left(\tilde{\bm{\lambda}}_{j}\mid \bm{\lambda}_{j}\right)p\left(\bm{\lambda}_{j}\right)$ and the marginal integral,
\begin{equation}
    p\left(\tilde{\bm{\lambda}}_{j}\right)=\int \mathcal{N}\left( \tilde{\bm{\lambda}}_{j};\,\bm{\lambda}_j,\mathbf{F}_{j}^{-1}\right)\mathcal{N}\left(\bm{\lambda}_{j};\,\bm{0},\bm{\Sigma}\right)\,d\bm{\lambda}_{j},
    \label{eq:gaussian-mix}
\end{equation}
% FIX: the first factor was written N(lambda_j; 0, F_j^{-1}) -- an expression that does not even mention lambda-tilde_j on the right-hand side, so it could not be a valid derivation of p(lambda-tilde_j). Given Equation~\eqref{eq:observation_model} directly above (epsilon_j = lambda-tilde_j - lambda_j ~ N(0,F_j^{-1}), i.e. lambda-tilde_j | lambda_j ~ N(lambda_j, F_j^{-1})) and the correct final result in Equation~\eqref{eq:fh-marginal-likelihood} just below, the first factor here must be N(lambda-tilde_j; lambda_j, F_j^{-1}); changed accordingly.
which is an approximation. Strictly speaking, $\mathbf{F}_j$ depends on the realized $\bm{\lambda}_j$ through the multinomial information matrix, so Equation~\eqref{eq:observation_model} is a conditional statement, in which $\mathbf{F}_{j}$ is a function of $\bm{\lambda}_{j}$.
% FIX: "F_j is the function of lambda_j" -> "F_j is a function of lambda_j".
Thus, the integral~\eqref{eq:gaussian-mix} is a Gaussian scale mixture, rather than an exact Gaussian.
% FIX: "an Gaussian scale mixture" -> "a Gaussian scale mixture".
We write it as an approximated formula,
\begin{equation} 
    l_{\mathrm{FH}}(\bm{\Sigma}) \approx \sum_{j=1}^{J} \log \mathcal{N}\left(\tilde{\bm{\lambda}}_j;\, \mathbf{0}, \bm{\Sigma} + \mathbf{F}_j^{-1}\right). 
    \label{eq:fh-marginal-likelihood} 
\end{equation}
The right side is exactly the marginal likelihood of the multivariate Fay-Herriot model~\citep{Herriot1979, Pfeffermann2011}. $\tilde{\bm{\lambda}}_j$ is a direct estimate with known sampling covariance $\mathbf{F}_j^{-1}$, while $\bm{\Sigma}$ is the between-group covariance. When all $\mathbf{F}_j$ are equal, Equation~\eqref{eq:fh-marginal-likelihood} is exactly the MLE of $l_{\text{FH}}(\bm{\Sigma})$. When $\mathbf{F}_j$ are not equal, which naturally arises since the information from different groups can differ and our model allows different $I_j$, the debias estimator in Equation~\eqref{eq:debias-Sigma} remains consistent, and each term $\tilde{\bm{\lambda}}_j\tilde{\bm{\lambda}}_j^{\top}-\mathbf{F}_j^{-1}$ has expectation $\bm{\Sigma}$ regardless of $\mathbf{F}_j$, so the unweighted average converges to $\bm{\Sigma}$ as $J\to\infty$. It is not, however, fully efficient. Maximizing $l_{\text{FH}}(\bm{\Sigma})$ instead weights each group by its precision $(\bm{\Sigma}+\mathbf{F}_j^{-1})^{-1}$, so groups with higher information contribute more to the estimate of $\bm{\Sigma}$. Equal weighting, as used by debias, discards this information and results in a larger asymptotic variance.

\paragraph*{E-step.}
\label{para:fhem-estep}
Equation~\eqref{eq:fh-marginal-likelihood} carries the approximation error discussed above. To avoid it, we introduce a genuine EM algorithm, where the missing data are $\bm{\lambda}_j$ and the observations are $\tilde{\bm{\lambda}}_j$. Given the current $\bm{\Sigma}$ and denote $\mathbf{V}_j = \left(\mathbf{F}_j + \bm{\Sigma}^{-1}\right)^{-1}$ as the covariance matrix of $\bm{\lambda}_{j}\mid \tilde{\bm{\lambda}}_{j}$, the joint log-likelihood is,
\begin{align}
    \log p\left(\bm{\lambda}_{j},\tilde{\bm{\lambda}}_{j}\right)&=-\frac{1}{2}\left(\bm{\lambda}_{j}-\tilde{\bm{\lambda}}_{j}\right)^{\top}\mathbf{F}_{j}\left(\bm{\lambda}_{j}-\tilde{\bm{\lambda}}_{j}\right)-\frac{1}{2}\bm{\lambda}_{j}^{\top}\bm{\Sigma}^{-1}\bm{\lambda}_{j}\\
    &=-\frac{1}{2}\left[\bm{\lambda}_{j}^{\top}\left(\mathbf{F}_{j}+\bm{\Sigma}^{-1}\right)\bm{\lambda}_{j}-2\tilde{\bm{\lambda}}_{j}^{\top}\mathbf{F}_{j}\bm{\lambda}_{j}\right]\\
    &=-\frac{1}{2}\left[\left(\bm{\lambda}_{j} - \mathbf{V}_{j} \mathbf{F}_{j} \tilde{\bm{\lambda}}_{j}\right)^{\top} \mathbf{V}_{j}^{-1} \left(\bm{\lambda}_{j} - \mathbf{V}_{j} \mathbf{F}_{j} \tilde{\bm{\lambda}}_{j}\right) - \tilde{\bm{\lambda}}_{j}^{\top} \mathbf{F}_{j} \mathbf{V}_{j} \mathbf{F}_{j} \tilde{\bm{\lambda}}_{j}\right],
    \label{eq:joint-log-density}
\end{align}
up to a constant independent of $\bm{\lambda}_{j}$.
% FIX: "independent to" -> "independent of".
Notice that $\bm{\lambda}_{j}$ only appears in the first term, the posterior therefore is given by $\mathbf{V}_{j}$ and the posterior mean $\mathbf{m}_{j} = \mathbf{V}_{j} \mathbf{F}_{j} \tilde{\bm{\lambda}}_{j}$,
\begin{equation}
    \bm{\lambda}_{j}\mid \tilde{\bm{\lambda}}_{j}\sim \mathcal{N}\left(\mathbf{m}_{j},\mathbf{V}_{j}\right).
    \label{eq:fh-posterior}
\end{equation}
Interestingly, we note that the posterior mean $\mathbf{m}_{j} = \mathbf{V}_{j} \mathbf{F}_{j} \tilde{\bm{\lambda}}_{j}=\left(\mathbf{F}_{j}+\bm{\Sigma}^{-1}\right)^{-1}\mathbf{F}_{j}\tilde{\bm{\lambda}}_{j}$ has the same form as the shrinkage factor in Equation~\eqref{eq:classical-bayes2}, confirming that $\mathbf{m}_{j}$ is the posterior mean of $\bm{\lambda}_{j}$ given $\bm{\Sigma}$.
% FIX: "we extend the posterior mean..., which is the similar formula as... and implies m_j to be the posterior mean" -- reworded for grammar, same claim.

We call back the structure of standard EM-algorithm and find the complete log-likelihood. Since $\bm{\Sigma}$ is not contained in $\tilde{\bm{\lambda}}_{j}\mid \bm{\lambda}_{j}$, we have,
\begin{align}
    l_{\text{complete}}\left(\bm{\Sigma}\right)&=\sum_{j}\log \mathcal{N}\left(\bm{\lambda}_{j};\,\bm{0},\bm{\Sigma}\right)=-\frac{J}{2}\log|\bm{\Sigma}| - \frac{1}{2}\sum_{j}\bm{\lambda}_j^{\top}\bm{\Sigma}^{-1}\bm{\lambda}_j + \text{const}\\
    &= -\frac{J}{2}\log|\bm{\Sigma}| - \frac{1}{2}\operatorname{tr}\left(\bm{\Sigma}^{-1}\sum_{j}\bm{\lambda}_j\bm{\lambda}_j^{\top}\right) + \text{const}
    \label{eq:fhem-l-complete}
\end{align}
Applying $\mathbb{E}\left[\mathbf{X}\mathbf{X}^{\top}\right]=\operatorname{Var}\left(\mathbf{X}\right)+\mathbb{E}\left[\mathbf{X}\right]\mathbb{E}\left[\mathbf{X}\right]^{\top}$, we substitute $\mathbb{E}\left[\bm{\lambda}_{j}\bm{\lambda}_{j}^{\top}\mid \tilde{\bm{\lambda}}_{j}\right]=\mathbf{V}_{j}+\mathbf{m}_{j}\mathbf{m}_{j}^{\top}$ into~\eqref{eq:fhem-l-complete} and obtain the Q-function for maximization,
\begin{equation}
    Q(\bm{\Sigma} \mid \bm{\Sigma}^{\text{old}}) = -\frac{J}{2}\log|\bm{\Sigma}| - \frac{1}{2}\operatorname{tr}\left(\bm{\Sigma}^{-1}\sum_{j}\left(\mathbf{m}_j\mathbf{m}_j^{\top} + \mathbf{V}_j\right)\right) + \text{const}.
    \label{eq:fhem-Q-function}
\end{equation}

\paragraph*{M-step.}
\label{para:fhem-mstep}
Calculate the derivative of~\eqref{eq:fhem-Q-function} and obtain the solution,
\begin{equation} \mathbf{S} = \frac{1}{J} \sum_{j=1}^{J} \left( \mathbf{m}_j \mathbf{m}_j^{\top} + \mathbf{V}_j \right). 
\label{eq:fhem_mstep} 
\end{equation}
As in the standard Laplace-EM algorithm in Section~\ref{sec: EM-wb}, we apply the closed-form PPCA solution in Equations~\eqref{eq:tipping-ppca_b}-\eqref{eq:tipping-ppca_sigma} to $\mathbf{S}$ to decompose $\mathbf{B}$ and $\sigma^{2}$.
% FIX: "As the similar as standard Laplace-EM algorithm..." -> "As in the standard Laplace-EM algorithm...". Also note: this reuses Equations~\eqref{eq:tipping-ppca_b}-\eqref{eq:tipping-ppca_sigma}, so the Q-K vs Q-1-K flag above applies here too -- fixing that equation once fixes both places it's used.

FH-EM algorithm can solve the Heywood problem that happens in the debias method (Section~\ref{subsubsec:debias}) by establishing a strictly positive definite matrix $\mathbf{S}$. The collapse mechanism is thus structurally blocked rather than simply not appearing in a particular experiment. This is the specific reason why the ``Heywood'' problem is structurally corrected by the new method.

We tested FH-EM algorithm on standard EM-algorithm setting without Woodbury formula usage. Since these optimized processes work well on standard EM algorithm, we can also use the Woodbury to reduce the calculation cost and complexity, which is introduced in the next section. 

\paragraph*{Algorithm.}
\label{para:fhem-algorithm}
Algorithm~\ref{alg:fhem} summarizes the FH-EM procedure derived above. The inputs $\tilde{\bm{\lambda}}_j$ and $\mathbf{F}_j$ come from Stage~1 and remain fixed throughout. Each iteration alternates between the E-step in Equations~\eqref{eq:fh-posterior} and the M-step in Equation~\eqref{eq:fhem_mstep}, followed by a closed-form PPCA update. Both terms in the M-step are positive semi-definite by construction, so $\mathbf{S}$ is always positive definite and the collapse described in Level~2 of Table~\ref{tab:heywood-levels} cannot occur.
% NOTE (not a change, just a heads-up): the referenced Table~\ref{tab:heywood-levels} currently sits inside an \iffalse ... \fi block below, so it will not appear in the compiled PDF as the thesis stands. Flagging in case that was accidental.
Convergence is monitored through the marginal log-likelihood $l_{\text{FH}}(\bm{\Sigma})$ rather than the change in $\bm{\Sigma}$ itself, since the latter can declare convergence prematurely.
% FLAG (not changed): elsewhere in this thesis the decided convergence rule for the Woodbury EM was a JOINT absolute-tolerance check on both the log-evidence change and the sigma^2 change, specifically because your advisor wanted absolute (not relative) tolerance and because log-evidence alone was shown to be non-monotone in some of your earlier traces. This sentence recommends log-likelihood alone (explicitly rejecting monitoring Sigma's own change) for FH-EM. If FH-EM's log-likelihood genuinely is monotone (per (A5)/(A7) above, it should be, for the inner sub-problem), a log-likelihood-only rule may be fine here and simply different from the Woodbury-EM rule for a good reason -- but please confirm this is deliberate rather than an inconsistency with the joint-criterion decision elsewhere.

\begin{algorithm}[htbp]
\caption{FH-EM: maximum-likelihood estimation of $\bm{\Sigma}$}
\label{alg:fhem}
\begin{algorithmic}[1]
\Require Per-group pure-data MLEs $\tilde{\bm{\lambda}}_j$ and Fisher information matrices $\mathbf{F}_j$, $j=1,\dots,J$, obtained from Stage~1 (isotropic-prior Laplace approximation)
\Require Initial values $\mathbf{B}^{(0)}$, $\sigma^{2(0)}$; tolerance $\mathrm{tol}$; maximum iterations $T$
\Ensure $\hat{\mathbf{B}}$, $\hat\sigma^2$
\State $\bm{\Sigma}^{(0)} \gets \mathbf{B}^{(0)}(\mathbf{B}^{(0)})^{\top} + \sigma^{2(0)}\mathbf{I}_Q$
\State $l^{(0)} \gets l_{\text{FH}}(\bm{\Sigma}^{(0)})$
\For{$t = 0, 1, 2, \dots, T-1$}
    \Statex \hspace{\algorithmicindent}\textit{E-step}
    \For{$j = 1, \dots, J$}
        \State $\mathbf{V}_j \gets \left(\mathbf{F}_j + (\bm{\Sigma}^{(t)})^{-1}\right)^{-1}$
        \State $\mathbf{m}_j \gets \mathbf{V}_j \mathbf{F}_j \tilde{\bm{\lambda}}_j$
    \EndFor
    \Statex \hspace{\algorithmicindent}\textit{M-step}
    \State $\mathbf{S} \gets \dfrac{1}{J}\displaystyle\sum_{j=1}^{J}\left(\mathbf{m}_j\mathbf{m}_j^{\top} + \mathbf{V}_j\right)$
    \State $\left(\mathbf{B}^{(t+1)}, \sigma^{2(t+1)}\right) \gets \mathrm{PPCA}(\mathbf{S}, K)$ \Comment{closed form, Eqs.~\eqref{eq:tipping-ppca_b}-\eqref{eq:tipping-ppca_sigma}}
    \State $\bm{\Sigma}^{(t+1)} \gets \mathbf{B}^{(t+1)}(\mathbf{B}^{(t+1)})^{\top} + \sigma^{2(t+1)}\mathbf{I}_Q$
    \State $l^{(t+1)} \gets l_{\text{FH}}(\bm{\Sigma}^{(t+1)})$
    \If{$\left|l^{(t+1)} - l^{(t)}\right| < \mathrm{tol}$}
    \State \textbf{break}
    \EndIf
\EndFor
\State \Return $\hat{\mathbf{B}} \gets \mathbf{B}^{(t+1)}$, $\hat\sigma^2 \gets \sigma^{2(t+1)}$
\end{algorithmic}
\end{algorithm}

\iffalse
\subsection{Solutions}
\label{subsec:solutions}

We have used the term Heywood problem loosely so far. It actually refers to three distinct failure mechanisms in this thesis, not one. Table~\ref{tab:heywood-levels} separates them and lists which method addresses each one.

\begin{table}[htbp]
\centering
\begin{tabular}{p{1.6cm}p{6.3cm}p{1.6cm}p{4.3cm}}
\toprule
Level & Issue Description & Fixed by & Fixed method \\
\midrule
Issue 1 & 
The E-step prior~\eqref{eq: Overall notation structure} depends on the current $\hat{\mathbf{B}}$, causing $\hat{\bm{\lambda}}_j$ to shrink less along the directions $\hat{\mathbf{B}}$ has already identified. The M-step then enhance the self-confirming loop. &
\textsc{iso} &
Replace the E-step prior with isotropic working prior $\tau^2\mathbf{I}$. The shrinkage weight no longer depends on the direction of $\hat{\mathbf{B}}$, so the loop is broken. \\
\addlinespace
Issue 2 &
The debias estimator~\eqref{eq:debias-Sigma} subtracts one positive semi-definite matrix from another. Nothing prevents the result from leaving the positive semi-definite cone, which drives $\hat\sigma^2$ to collapse. &
\textsc{fh-em} &
Matrix $\mathbf{S}$ in M-step~\eqref{eq:fhem_mstep} is positive definite by construction, so this collapse cannot occur. \\
\addlinespace
Issue 3 &
When the true $\sigma^2$ is small or $K$ is misspecified, the likelihood itself favors a solution at or near the boundary. A property of the model and the data, not an artifact of the algorithm. &
None &
Not addressed by \textsc{iso}, debias or \textsc{fh-em}. Would require an adjusted likelihood penalty in the M-step, which is left as future work. \\
\bottomrule
\end{tabular}
\caption{Three distinct mechanisms behind the Heywood problem in this thesis, and which correction method addresses each one.}
\label{tab:heywood-levels}
\end{table}
\fi

\section{Consistency \& Asymptotics}
\label{sec:asymptotical-properties}